\setlength{\absparsep}{18pt} 

\begin{resumo}[Resumo]
 \begin{otherlanguage*}{brazil}
   
   \vspace{\onelineskip}
   \noindent 
   
O uso frequente de aplicações de inteligência artificial por estudantes da Educação Básica não garante que eles entendam como dados e regras orientam o funcionamento dessas ferramentas no dia a dia. Este projeto de qualificação parte desse problema. Com o propósito de analisar como os estudantes explicam o modo pelo qual modelos computacionais partem de dados para produzir previsões numéricas, separar classes ou formar grupos, busca-se desenvolver e analisar uma sequência didática apoiada por uma plataforma interativa para o ensino introdutório de \textit{machine learning}, como conteúdo de Computação no Ensino Médio. Com referência no pensamento computacional e nas orientações da Resolução CNE/CEB nº 1/2022, a proposta recorre ao uso de múltiplas representações, com as aulas indo da previsão inicial ao experimento prático, de modo a dar base à observação e à explicação. O produto educacional reúne oito encontros. Cobrindo regressão linear, máquina de vetores de suporte (\textit{SVM}), \textit{k}-vizinhos mais próximos (\textit{k}-NN), árvore de decisão e \textit{k}-means, essa sequência faz uso de uma plataforma web executada localmente, com dados sintéticos tabulares. A pesquisa tem natureza aplicada e abordagem qualitativa. Apoiado por dados quantitativos descritivos e conduzido como intervenção pedagógica numa turma do próprio pesquisador, o estudo reunirá em seu corpus o diagnóstico inicial, as notas docentes, três registros intermediários e o instrumento final com seu desafio de transferência. Uma rubrica guiará o exame dos registros. A partir desse instrumento, será identificado o que muda ou permanece nas dificuldades dos estudantes ao explicarem como os dados e a finalidade da tarefa se relacionam com regras, parâmetros e resultados, considerando-se ainda o reconhecimento dos limites das decisões automatizadas. Frequências e transições entre níveis descreverão o conjunto, sem atribuir à plataforma um efeito causal isolado.
   
   \textbf{Palavras-chave}: Ensino de Computação; Machine Learning; Ensino Médio; Sequência Didática; Plataforma Interativa; Produto Educacional.

 \end{otherlanguage*}
\end{resumo}

% O resumo em inglês (Abstract) será incluído após a validação da versão portuguesa.
